\subsection{The seven-invertible Morita candidates}
\label{subsec:focus-r3472-candidates}

Put $r=\sqrt2$, $\delta=7+5r$, and $D=700+490r$.
The following statement concerns the seven-invertible alternative in the
preceding Morita restriction. It does not assert that this alternative occurs.

\begin{proposition}\label{prop:focus-r3472-two-rings}
Every seven-invertible Morita partner of a categorification of
$\sha{R8_3_3472}$ has one of the following two based rings. Their basis is
$\{g_i,X_i:i\in\mathbb Z/7\mathbb Z\}$, with
\[
 g_i g_j=g_{i+j},\qquad g_iX_j=X_{i+j},\qquad
 X_i g_j=X_{i+aj},\qquad
 X_iX_j=g_{i+aj}+2\sum_kX_k,
 \quad a\in\{1,-1\}.
\]
Both displayed multiplications are associative fusion rings. No associator
or Morita identification with the source ring is established here.
\end{proposition}
\begin{proof}
The preceding dimension restriction gives $1^7,\delta^7$.
Since $\delta^2=1+14\delta$, reciprocity makes both stabilizers of every
noninvertible trivial. Write $X_i=g^iX$ and $Xg=g^aX$, where $a$ is a
unit modulo seven. The unique invertible summand of $X^2$ is $g^b$.
Comparing the invertible parts in $X^2g=g^{a^2}X^2$ gives $a^2=1$.
Reciprocity gives $X^*=g^{-b}X$, and applying duality twice gives
$(a-1)b=0$. Thus $b=0$ if $a=-1$, and replacing $X$ by $g^cX$ makes
$b=0$ if $a=1$. In particular $X$ is self-dual.

For $a=-1$, the identity $X^2g=gX^2$ makes the noninvertible
coefficients invariant under translation by two. They sum to fourteen,
so all are two. For $a=1$, write
$X^2=1+m(g)X$ with $m_j=m_{-j}\geq0$ and $\sum_jm_j=14$.
Pass to the pseudo-unitary Galois conjugate supplied by the source
analysis; this does not change any fusion coefficient. All simple
dimensions in its common center belong to $\mathbb Q(r)$, by forgetting
to the source category. The formal-codegree induction formula
$\operatorname{FPdim}(A_E)=D/f_E$ therefore gives
$f_E\in\mathbb Q(r)$ for every irreducible representation $E$.

Choose a primitive seventh root $\zeta$. The characters above
$g\mapsto\zeta$ have $X\mapsto x$, where
$x^2-tx-1=0$ and $t=m(\zeta)$ is real. Their formal codegrees are
$7(1+x^2)$. Consequently
\[
 t^2=\frac{(x^2-1)^2}{x^2}
 \in\mathbb Q(r)\cap\mathbb Q(\zeta+\zeta^{-1})=\mathbb Q.
\]
The latter field has odd degree three, hence $t\in\mathbb Q$.
The polynomial $\Phi_7$ now forces $m_0=A$ and $m_j=B$ for $j\ne0$,
with $A+6B=14$. The possibilities are $(A,B)=(14,0),(8,1),(2,2)$.
The first contains the impossible rank-two categorification
$X^2=1+14X$. The second has formal codegrees
$(371\pm49\sqrt{53})/2$, outside $\mathbb Q(r)$. Hence $A=B=2$.
The complete tensors, dualities, dimensions, reciprocity identities and
all $14^4$ associativity equations for each ring are checked in
\path{evidence/r3472_focus_20261022/constructions/verify_constructions.py}.
\end{proof}

\subsection{An obstruction to all direct Radchenko--Wheeler parent routes}
\label{subsec:focus-r3472-rw}

We use the explicit near-group categories of Radchenko--Wheeler,
\emph{Real quadratic fields and finite quantum dilogarithms I},
arXiv:2609.21892v2, Sections~4.3 and~5.4. The scope below is their
construction from a matrix $\gamma\in\mathrm{SL}_2(\mathbb Z)$;
no classification of all categorifications of the parent fusion ring
is asserted.

\begin{lemma}\label{lem:focus-rw-halfbraiding}
For a Radchenko--Wheeler matrix input with
$\operatorname{tr}(\gamma)=N+2$, put
\[
 G=\mathbb Z^2/\mathbb Z^2(\gamma-I),\qquad
 \chi(u,v)=\exp\!\left(\frac{2\pi i}{N}
                  \omega(u(\gamma-I),v)\right).
\]
The central invertibles embed into $G$, and their self-braidings agree
with $q(g)=\chi(g,g)^{-1}$ on their image.
\end{lemma}
\begin{proof}
The strict algebraic construction gives
$\alpha_h\rho=\rho$,
$\rho\alpha_g=\operatorname{Ad}(U_g)\rho$, and
$\alpha_h(U_g)=\chi(h,g)U_g$.
The half-braiding of a central lift of $\alpha_g$ with $\rho$ is a
nonzero scalar times $U_g$, since $\rho$ is simple. Its hexagon at
$\alpha_h\rho=\rho$ forces its half-braiding with $\alpha_h$ to be
$\chi(h,g)^{-1}$. This uses neither a star operation nor unitarity.
The forgetful kernel on central invertibles consists of tensor
automorphisms of the identity. It is trivial because the universal
grading is trivial: $\rho$ and every $\alpha_g$ occur in $\rho^2$.
\end{proof}

\begin{lemma}\label{lem:focus-rw-odd-square}
For $N=4s^2$ with $s$ odd, the bosonic part of the radical of
$(G,q)$ has order at most two.
\end{lemma}
\begin{proof}
The polarization is $\chi^{-2}$, with radical $G[2]$.
The two-primary part of the two-generated group $G$ has order four.
If it is cyclic, $|G[2]|=2$. Otherwise
$\gamma-I=2A$, where
\[
 \operatorname{tr}A=2s^2\equiv2\pmod4,\qquad
 \det A=-s^2\equiv3\pmod4.
\]
The two rows of $A$ represent a basis of $G[2]$. For
$A=\left(\begin{smallmatrix}a&b\\c&d\end{smallmatrix}\right)$,
the corresponding diagonal bicharacter values are $(-1)^b$ and
$(-1)^c$. Triviality of $q|_{G[2]}$ would make $b,c$ even.
Then $a,d$ are odd and their sum is two modulo four, so $ad$ is one
modulo four. This contradicts $\det A\equiv3\pmod4$.
Thus $q$ is a nontrivial character on the four-element radical,
and its kernel has order two.
\end{proof}

\begin{theorem}\label{thm:focus-rw-pointed-stabilization}
Let $s$ be odd and positive, set
$\delta_s=s+\sqrt{s^2+1}$ and $D_s=s(1+\delta_s^2)$.
No Radchenko--Wheeler parent with group order $4s^2$, tensored with
an arbitrary pointed fusion category, admits an ordinary
de-equivariantization of dimension $D_s$.
In particular, no such direct route, including noncyclic parents and
pointed tensor factors, supplies a Morita representative for
$\sha{R8_3_3472}$.
\end{theorem}
\begin{proof}
Let the pointed factor $\mathcal P$ have dimension $m$.
The parent has dimension $4sD_s$, so the required Tannakian group
would have order $h=4sm$. Its fully faithful forgetful image is pointed:
simple representations have integer dimensions, whereas all
noninvertible dimensions in the parent tensor product are irrational.
Thus its simples form an isotropic subgroup $H$ of order $h$ in the
central invertibles.

Put $L=\operatorname{Inv}(Z(\mathcal P))$, $k=|L|$, and let the
radical of its self-braiding form have order $t$.
The centralizer dimension formula in the nondegenerate category
$Z(\mathcal P)$ gives $kt\leq m^2$.
Lemma~\ref{lem:focus-rw-halfbraiding} embeds $H$, preserving its
quadratic form, into $G\times L$. Its intersection with the radical
has order at most $2t$: in the cyclic two-primary case the full
radical has that order, and in the other case the radical quadratic
form is nontrivial already on $G[2]$, so its kernel has that order.
For any isotropic subgroup of a finite premetric group $A$,
\[
 |H|^2\leq |A|\,|H\cap\operatorname{rad}(A)|,
\]
because $H\subseteq H^\perp$ and
$|H^\perp|=|A|\,|H\cap\operatorname{rad}(A)|/|H|$.
It follows that
$h^2\leq4s^2k(2t)\leq8s^2m^2$, contradicting
$h^2=16s^2m^2$.

For the stated target, no other near-group parent order can have
the required rational dimension scale. The global dimension $E_N$
of $\mathrm{NG}(G,N)$, $|G|=N$, has trace $N(N+4)$ and norm
$N^2(N+4)$. Hence its scale-invariant trace-squared/norm ratio is
$N+4$. For $D$ the ratio is $1400^2/9800=200$, forcing $N=196$.
The case $s=7$ above therefore covers all Radchenko--Wheeler matrix
parents with a possible dimension match.
\end{proof}

This rules out the specified direct ordinary quotients, also followed by
a Morita equivalence. It does not address a prior Morita change of the
parent, arbitrary near-group categorifications with different input data,
generalized-Haagerup parents, or arbitrary chains of constructions.
The two rank-fourteen tensors remain unconstructed candidate based rings;
$\sha{R8_3_3472}$ remains unresolved.
